\chapter{DISCUSIÓN} \label{chap:discusion}

Los resultados experimentales obtenidos en este trabajo muestran que los modelos basados
en redes neuronales convolucionales con aprendizaje por transferencia constituyen una
alternativa viable para la detección automatizada de retinopatía diabética en imágenes
del fondo de ojo. A continuación se interpretan los hallazgos más relevantes,
contrastándolos con el marco teórico y con investigaciones previas en el área.

\section{Desempeño del clasificador binario y comparación con la literatura}

El modelo MobileNetV2 con filtrado y umbral óptimo alcanzó una sensibilidad de 96.5\% y
un F1-score de 0.96 en el conjunto de prueba interno, métricas que lo posicionan
favorablemente frente a trabajos previos. Gulshan et al. reportaron una sensibilidad del
97.5\% utilizando arquitecturas más profundas entrenadas con conjuntos de datos
considerablemente mayores \myfootcite{gulshan2016development}, lo cual sugiere que la
estrategia de preprocesamiento implementada en este trabajo logró compensar parcialmente
las limitaciones en volumen de datos mediante la mejora de la calidad de las
características visuales extraídas.

Es importante precisar el alcance de las siguientes explicaciones: este trabajo compara
arquitecturas completas entre sí y no realiza experimentos de ablación que aíslen la
contribución individual de cada componente arquitectónico al desempeño observado. Por
ello, las relaciones entre características de diseño y desempeño que se plantean a
continuación deben entenderse como interpretaciones posibles, sustentadas en literatura
previa, y no como causas demostradas experimentalmente por este trabajo.

Una interpretación posible de la superioridad de MobileNetV2 sobre las demás
arquitecturas evaluadas (VGG19, ResNet50V2 e InceptionV3) se relaciona con
características propias de su diseño. Los bloques residuales invertidos de MobileNetV2,
que emplean convoluciones separables en profundidad junto con cuellos de botella
lineales, se han asociado con una forma de regularización implícita que podría resultar
beneficiosa cuando el conjunto de datos disponible es limitado
\myfootcite{sandler2018mobilenetv2}, lo que ofrece una posible explicación de por qué
este modelo mostró menor propensión al sobreajuste que arquitecturas más grandes en este
experimento. De manera similar, se ha planteado que la ausencia de activaciones no
lineales en las capas de proyección ayuda a preservar información de baja
dimensionalidad \myfootcite{sandler2018mobilenetv2}, lo cual podría ser relevante para
capturar características sutiles como las lesiones microvasculares en etapas tempranas de
la RD, aunque este trabajo no verificó dicho mecanismo de forma directa.

Los resultados contrastan con el desempeño de VGG19, que a pesar de su profundidad y
capacidad de extracción de características, mostró métricas inferiores (F1-score de 0.86
sin filtrado). Una posible explicación es que esta arquitectura, caracterizada por un
elevado número de parámetros en sus capas completamente conectadas y la ausencia de
conexiones residuales, sea más propensa al sobreajuste en este escenario de datos
limitados; sin embargo, este trabajo no midió directamente el grado de sobreajuste de
cada modelo (por ejemplo, mediante la brecha entre error de entrenamiento y validación),
por lo que esta explicación permanece como hipótesis. Es consistente, en todo caso, con
lo reportado por Kandel y Castelli, quienes identificaron que arquitecturas más compactas
tienden a generalizar mejor en dominios médicos con datos limitados
\myfootcite{kandel2020transfer}.

\section{Impacto del preprocesamiento en el rendimiento del modelo}

La implementación del algoritmo de filtrado, que incluyó normalización del color mediante
sustracción del promedio local con filtro gaussiano y recorte del fondo oscuro de la
imagen, produjo mejoras consistentes en todas las arquitecturas evaluadas. El incremento
del F1-score de 0.95 a 0.96 en MobileNetV2, aunque aparentemente modesto, se acompañó de
una reducción de falsos positivos (de 6 a 4 casos) y de falsos negativos (de 5 a 3 casos),
con implicaciones clínicas relevantes al disminuir tanto las derivaciones innecesarias
como los casos de RD no detectados en un escenario de tamizaje.

Este resultado es consistente con la importancia del preprocesamiento destacada en la
revisión de literatura: Tabtaba y Ata reportaron mejoras similares al aplicar técnicas de
ecualización adaptativa del histograma sobre imágenes de fondo de ojo
\myfootcite{tabtaba2024diabetic}. La normalización de color implementada en este trabajo
busca mitigar las variaciones de brillo y contraste inherentes a la captura de imágenes
en entornos clínicos heterogéneos, donde las condiciones de iluminación y los equipos de
adquisición varían considerablemente; el recorte del fondo oscuro, por su parte, elimina
regiones no informativas que podrían introducir ruido en el proceso de extracción de
características.

\section{Validación externa y capacidad de generalización}

La evaluación del modelo en conjuntos de datos externos (FOSCAL, IDRiD y MESSIDOR-2)
mostró un desempeño variable según el conjunto, con una disminución en las métricas
respecto al conjunto de prueba interno de Kaggle (Sección~\ref{subsec:validacion-externa}).
Este hallazgo —una caída de desempeño al evaluar fuera del dominio de entrenamiento— no
es particular de este trabajo: es un fenómeno documentado en la literatura reciente sobre
generalización cruzada de modelos de clasificación de RD. Hashmi reporta, por ejemplo,
una caída de sensibilidad para RD referible de 0.879 en el dominio de entrenamiento
(APTOS) a 0.620 al evaluar sobre MESSIDOR-2 sin reentrenamiento, señalando que las
métricas en dominio tienden a sobreestimar el desempeño real desplegable
\myfootcite{hashmi2026crossdomain}. Este patrón resulta relevante para interpretar por
qué el desempeño de este trabajo también varía entre datasets externos.

Las diferencias entre FOSCAL, IDRiD y MESSIDOR-2 no se limitan a la calidad visual de las
imágenes: involucran distinto modelo de cámara, resolución, protocolo de captura y
población de origen (india, francesa y colombiana, respectivamente; véase la
Sección~\ref{sec:conjuntos-datos}), factores que la literatura sobre cambio de dominio
identifica como determinantes de la transferibilidad de un modelo entrenado en un dominio
distinto \myfootcite{hashmi2026crossdomain}. Dado que estas variables no fueron
controladas ni desagregadas en el análisis, no es posible atribuir con precisión qué
proporción de la variación observada corresponde a cada factor; esto se reconoce como
limitación en la Sección~\ref{sec:limitaciones-discusion}.

En conjunto, estos resultados sugieren que el modelo mantiene un desempeño razonable
fuera del dominio de entrenamiento, comparable en magnitud a la caída reportada en
trabajos similares, pero no permiten afirmar una capacidad de generalización robusta o
universal: como en el caso documentado por Hashmi, las métricas del conjunto de
entrenamiento tienden a sobreestimar el desempeño esperable en un despliegue real, y una
conclusión más definitiva requeriría aislar el efecto de cada fuente de variación entre
datasets \myfootcite{hashmi2026crossdomain}.

\section{Análisis del clasificador multiclase y el efecto del canal de color}

Los experimentos de clasificación multiclase revelaron hallazgos relevantes respecto al
procesamiento del color en imágenes de fondo de ojo. El canal verde demostró
consistentemente el mejor desempeño (exactitud de 0.55, sensibilidad promedio de 0.56),
superando tanto a la imagen RGB completa (+11.8\% en F1-score promedio de las cuatro
arquitecturas evaluadas) como a los canales rojo (+25.2\%) y azul (+26.0\%) de manera
individual (Tabla~\ref{tab:mejora-canal-verde}). Esta ventaja se concentra
particularmente en la clase Severa (RD3), donde la sensibilidad obtenida con el canal
verde supera en más de 85\% a la del canal rojo y en más de 50\% a la del canal azul en las dos
arquitecturas de mejor desempeño (Sección~\ref{subsec:analisis-canales}). Este resultado
tiene fundamento fisiológico: el canal verde proporciona mejor contraste para la
visualización de estructuras vasculares y lesiones hemorrágicas debido a la absorción
diferencial de la hemoglobina en esta longitud de onda, el mismo principio que sustenta
el uso clínico histórico de filtros verdes en oftalmoscopia para realzar hemorragias y
exudados \myfootcite{pao2020bichannel}. Los microaneurismas, hemorragias y anomalías
microvasculares intrarretinianas, signos que definen clínicamente la etapa severa de la
retinopatía diabética, presentan mayor visibilidad en este espectro cromático, lo que es
consistente con la ganancia de sensibilidad observada precisamente en esa clase. La ventaja no
es, sin embargo, absoluta: en la clase PDR (RD4), asociada a la neoformación vascular, el
canal rojo iguala o supera al verde, lo que sugiere que ningún canal individual captura
por sí solo toda la información relevante en las etapas más avanzadas de la enfermedad.

Este hallazgo es consistente con la literatura disponible sobre análisis de canales
individuales en RD, aunque dicha literatura resulta comparativamente escasa
(Sección~\ref{chap:trabajos-previos}). Pao et al. reportaron mejoras en la detección de
RD al extraer específicamente el canal verde para el cálculo de características de
entropía, superando el desempeño obtenido con la imagen en escala de grises
\myfootcite{pao2020bichannel}. Singh, Gupta y Dung, en un estudio que comparó canales
individuales de forma más directa, también encontraron diferencias de desempeño entre
canales, aunque limitado a una sola arquitectura y sin un análisis sistemático a través de
múltiples redes preentrenadas \myfootcite{singh2024finetuned}. El presente trabajo amplía
esta evidencia al comparar los cuatro canales (RGB, rojo, verde y azul) de forma
sistemática sobre cuatro arquitecturas distintas en un escenario multiclase, lo que
sustenta de manera más generalizable la ventaja del canal verde reportada previamente.

La dificultad persistente en la clasificación de las clases RD2 (moderada) y RD3 (severa)
observada en todos los experimentos constituye un hallazgo relevante. La sensibilidad
para RD2 fue consistentemente la más baja (valores entre 0.01 y 0.31 dependiendo de la
configuración), lo cual sugiere que las características distintivas de esta categoría no
quedan suficientemente representadas en las arquitecturas empleadas. Desde la perspectiva
clínica, la retinopatía diabética moderada representa una etapa de transición
caracterizada por la presencia de más microaneurismas que la etapa leve, pero sin
alcanzar los criterios que definen la etapa severa. Esta naturaleza transicional, con
límites difusos entre categorías adyacentes, explica en parte la confusión observada en
las matrices de confusión.

Este fenómeno de confusión entre clases intermedias ha sido documentado en trabajos
previos: Gayathri et al. reportaron dificultades similares en la discriminación de etapas
moderadas, proponiendo el uso de clasificadores en cascada o características multiescala
como estrategias de mitigación \myfootcite{gayathri2021diabetic}. Los resultados
obtenidos sugieren que futuras investigaciones podrían explorar técnicas específicas para
abordar este desafío, incluyendo funciones de pérdida focal que penalicen más severamente
los errores en clases difíciles, aumento de datos específico para clases intermedias, o
mecanismos de atención que enfoquen el análisis en regiones anatómicas específicas donde
las lesiones discriminantes tienden a manifestarse.

\section{Justificación de las decisiones arquitectónicas}

La configuración del clasificador, que incorpora bloques Multi-Head Attention posteriores
al extractor de características y capas GaussianNoise entre las capas densas, responde a
consideraciones específicas del dominio de aplicación. Los mecanismos de atención
permiten al modelo ponderar diferencialmente las regiones del mapa de características, lo
cual resulta particularmente relevante en imágenes de fondo de ojo donde las lesiones
ocupan fracciones pequeñas del área total de la imagen; esta capacidad de enfoque
selectivo, documentada en trabajos como el de Dosovitskiy et al.
\myfootcite{dosovitskiy2020image}, complementa la extracción de características locales
realizada por las capas convolucionales con información de contexto global.

La inyección de ruido gaussiano durante el entrenamiento actúa como regularizador
estocástico que expone al modelo a versiones perturbadas de las activaciones internas.
Esta estrategia, respaldada por los estudios de Zur et al., simula la variabilidad
inherente a la adquisición de imágenes médicas (diferencias en iluminación, ruido del
sensor, artefactos de preprocesamiento), forzando al clasificador a aprender
representaciones robustas que no dependan de valores exactos de activación
\myfootcite{zur2009noise}.

La selección de la función de activación ReLU para capas intermedias y Softmax para la
capa de salida sigue prácticas establecidas en la literatura. Asadi y Jiang demostraron
que esta combinación posee propiedades de aproximación universal, proporcionando
fundamento teórico a su amplia adopción empírica \myfootcite{asadi2020approximation}.
Además, la función Softmax garantiza que las salidas del clasificador multiclase
constituyan una distribución de probabilidad válida, facilitando la interpretación
clínica de las predicciones y la implementación de umbrales de confianza para casos que
requieran revisión por especialistas.

\section{Análisis de errores e implicaciones clínicas}

El análisis de las matrices de confusión revela patrones de error con implicaciones
clínicas diferenciadas. En el clasificador binario, la proporción de falsos negativos
(casos de RD no detectados) se mantuvo baja en las configuraciones con mejor desempeño
(3 a 5 casos sobre aproximadamente 105 por clase), lo cual es relevante en un contexto de
tamizaje donde la prioridad es maximizar la detección de casos que requieren atención
especializada. Los falsos positivos, aunque implican derivaciones innecesarias con su
correspondiente costo para el sistema de salud, no conllevan el riesgo de progresión de
la enfermedad asociado a los casos no detectados; esta asimetría de costos clínicos es la
que sustenta la priorización de la sensibilidad discutida en la
Sección~\ref{sec:metricas-evaluacion}.

En el clasificador multiclase, el patrón predominante de confusión entre clases
adyacentes es consistente con la naturaleza continua del proceso patológico subyacente.
La retinopatía diabética progresa de manera gradual, y los criterios de clasificación
establecidos imponen límites discretos sobre un espectro continuo de severidad. Esta
característica inherente al problema sugiere que métricas como el coeficiente kappa
ponderado, que consideran la proximidad entre clases predichas y reales, podrían
proporcionar una evaluación más apropiada del desempeño clínico del modelo que las
métricas de clasificación exacta utilizadas en este trabajo.

\section{Limitaciones del estudio y consideraciones metodológicas} \label{sec:limitaciones-discusion}

Es necesario reconocer varias limitaciones que contextualizan los resultados obtenidos.
El desbalance de clases presente en el conjunto de datos de Kaggle, donde la clase 0 (sin
retinopatía) representa aproximadamente el 73\% del total, constituye un factor que puede
sesgar el aprendizaje hacia la clase mayoritaria. Aunque se implementaron estrategias de
balanceo mediante submuestreo y aumento de datos, la distribución original refleja la
prevalencia real de la enfermedad en poblaciones de tamizaje, donde la mayoría de los
examinados no presenta la condición. Futuros trabajos podrían explorar técnicas
adicionales como el sobremuestreo sintético mediante redes generativas adversariales, el
uso de funciones de pérdida ponderadas por clase, o estrategias de aprendizaje con costos
asimétricos.

La variabilidad en la calidad de las imágenes del conjunto de datos, documentada en la
descripción original del repositorio de Kaggle, introduce ruido tanto en las imágenes
como en las etiquetas. La subjetividad inherente a la evaluación clínica, donde
diferentes especialistas pueden asignar grados distintos a la misma imagen, constituye
una fuente de incertidumbre que limita el techo de desempeño alcanzable por cualquier
modelo automatizado. Los intervalos de confianza obtenidos mediante bootstrap
(sensibilidad: 0.9510--1.0000; especificidad: 0.8511--0.9600) proporcionan una estimación
de la variabilidad de las métricas frente al conjunto de prueba utilizado, pero, como se
señaló en la Sección~\ref{subsec:umbral-decision}, no capturan la variabilidad asociada a
un reentrenamiento del modelo, a un cambio en la partición de los datos, ni la
incertidumbre asociada al etiquetado de referencia.

Adicionalmente, como se discutió en la sección de validación externa, las diferencias
entre FOSCAL, IDRiD y MESSIDOR-2 (cámara, resolución, protocolo de captura y población)
no fueron controladas ni analizadas de forma independiente, por lo que no es posible
determinar con precisión el peso relativo de cada factor sobre la disminución de
desempeño observada fuera del dominio de entrenamiento.

\section{Implicaciones para la práctica clínica}

Los resultados obtenidos sugieren que el modelo propuesto podría integrarse en programas
de tamizaje de retinopatía diabética como herramienta de apoyo al diagnóstico,
particularmente en contextos con acceso limitado a especialistas en oftalmología.
Considerando la situación epidemiológica en Colombia, donde se estima una proporción de
1.56 oftalmólogos por cada 100{,}000 habitantes con distribución geográfica desigual
\myfootcite{minsalud2018}, la implementación de sistemas de detección automatizada podría
contribuir a ampliar la cobertura del tamizaje y facilitar la detección temprana de casos
que requieren intervención.

No obstante, es fundamental enfatizar que estos sistemas deben concebirse como
herramientas de apoyo y no como sustitutos del criterio clínico especializado. El modelo
propuesto identifica casos sospechosos que requieren evaluación por un profesional de la
salud visual, pero la decisión diagnóstica final y la definición del plan terapéutico
deben permanecer bajo responsabilidad del especialista. La integración de estas
tecnologías en el flujo de trabajo clínico requiere, además, consideraciones adicionales
relacionadas con la explicabilidad de las predicciones, la gestión de la incertidumbre, y
la definición de protocolos claros para el manejo de casos discordantes entre el sistema
automatizado y la evaluación humana.
